\documentclass{ximera}

\begin{document}
	\author{Stitz-Zeager}
	\xmtitle{Exercises for the Inverse Trigonometric Functions}{}

\mfpicnumber{1} \opengraphsfile{ExercisesfortheInverseTrigonometricFunctions} % mfpic settings added 

\begin{question}
In Exercises \ref{exactvaluearcfirst} - \ref{exactvaluearclast}, find the exact value. (To enter $\pi$ in an answer box, type ``pi''.)

\begin{problem} \label{exactvaluearcfirst}
$\arcsin \left( -1 \right)$

\begin{solution}
$\arcsin \left( -1 \right) = -\frac{\pi}{2}$
\end{solution}
\end{problem}

\begin{problem}
$\arcsin \left( -\frac{\sqrt{3}}{2} \right)$

$\arcsin \left( -\frac{\sqrt{3}}{2} \right) = \answer{-\frac{\pi}{3}}$
\end{problem}

\begin{problem}
$\arcsin \left( -\frac{\sqrt{2}}{2} \right)$

\begin{solution}
$\arcsin \left( -\frac{\sqrt{2}}{2} \right) = -\frac{\pi}{4}$
\end{solution}
\end{problem}

\begin{problem}
$\arcsin \left( -\frac{1}{2} \right)$

$\arcsin \left( -\frac{1}{2} \right) = \answer{-\frac{\pi}{6}}$
\end{problem}

\begin{problem}
$\arcsin \left( 0 \right)$ 

\begin{solution}
$\arcsin \left( 0 \right) = 0$
\end{solution}
\end{problem}

\begin{problem}
$\arcsin \left( \frac{1}{2} \right)$

$\arcsin \left( \frac{1}{2} \right) = \answer{\frac{\pi}{6}}$
\end{problem}

\begin{problem}
$\arcsin \left( \frac{\sqrt{2}}{2} \right)$

\begin{solution}
$\arcsin \left( \frac{\sqrt{2}}{2} \right) = \frac{\pi}{4}$
\end{solution}
\end{problem}

\begin{problem}
$\arcsin \left( \frac{\sqrt{3}}{2} \right)$

$\arcsin \left( \frac{\sqrt{3}}{2} \right) = \answer{\frac{\pi}{3}}$
\end{problem}

\begin{problem}
$\arcsin \left( 1 \right)$

\begin{solution}
$\arcsin \left( 1 \right) = \frac{\pi}{2}$
\end{solution}
\end{problem}

\begin{problem}
 $\arccos \left( -1 \right)$

$\arccos \left( -1 \right) = \answer{\pi}$
\end{problem}
 
\begin{problem}
$\arccos \left( -\frac{\sqrt{3}}{2} \right)$

\begin{solution}
$\arccos \left( -\frac{\sqrt{3}}{2} \right) = \frac{5\pi}{6}$
\end{solution}
\end{problem}

\begin{problem}
$\arccos \left( -\frac{\sqrt{2}}{2} \right)$

$\arccos \left( -\frac{\sqrt{2}}{2} \right) = \answer{\frac{3\pi}{4}}$
\end{problem}

\begin{problem}
$\arccos \left( -\frac{1}{2} \right)$

\begin{solution}
$\arccos \left( -\frac{1}{2} \right) = \frac{2\pi}{3}$
\end{solution}
\end{problem}

\begin{problem}
$\arccos \left( 0 \right)$

$\arccos \left( 0 \right) = \answer{\frac{\pi}{2}}$
\end{problem}

\begin{problem}
$\arccos \left( \frac{1}{2} \right)$

\begin{solution}
$\arccos \left( \frac{1}{2} \right) = \frac{\pi}{3}$
\end{solution}
\end{problem}

\begin{problem}
$\arccos \left( \frac{\sqrt{2}}{2} \right)$

$\arccos \left( \frac{\sqrt{2}}{2} \right) = \answer{\frac{\pi}{4}}$
\end{problem}

\begin{problem}
$\arccos \left( \frac{\sqrt{3}}{2} \right)$

\begin{solution}
$\arccos \left( \frac{\sqrt{3}}{2} \right) = \frac{\pi}{6}$
\end{solution}
\end{problem}

\begin{problem}
$\arccos \left( 1 \right)$

$\arccos \left( 1 \right) = \answer{0}$
\end{problem}

\begin{problem}
$\arctan \left( -\sqrt{3} \right)$

\begin{solution}
$\arctan \left( -\sqrt{3} \right) = -\frac{\pi}{3}$
\end{solution}
\end{problem}

\begin{problem}
$\arctan \left( -1 \right)$

$\arctan \left( -1 \right) = \answer{-\frac{\pi}{4}}$
\end{problem}

\begin{problem}
$\arctan \left( -\frac{\sqrt{3}}{3} \right)$

\begin{solution}
$\arctan \left( -\frac{\sqrt{3}}{3} \right) = -\frac{\pi}{6}$
\end{solution}
\end{problem}

\begin{problem}
$\arctan \left( 0 \right)$

$\arctan \left( 0 \right) = \answer{0}$
\end{problem}

\begin{problem}
$\arctan \left( \frac{\sqrt{3}}{3} \right)$

\begin{solution}
$\arctan \left( \frac{\sqrt{3}}{3} \right) = \frac{\pi}{6}$
\end{solution}
\end{problem}

\begin{problem}
$\arctan \left( 1 \right)$

$\arctan \left( 1 \right) = \answer{\frac{\pi}{4}}$
\end{problem}


\begin{problem}
$\arctan \left( \sqrt{3} \right)$

\begin{solution}
$\arctan \left( \sqrt{3} \right) = \frac{\pi}{3}$
\end{solution}
\end{problem}

\begin{problem}
$\mbox{arccot} \left( -\sqrt{3} \right)$

$\mbox{arccot} \left( -\sqrt{3} \right) = \answer{\frac{5\pi}{6}}$
\end{problem}

\begin{problem}
$\mbox{arccot} \left( -1 \right)$

\begin{solution}
$\mbox{arccot} \left( -1 \right) = \frac{3\pi}{4}$
\end{solution}
\end{problem}

\begin{problem}
$\mbox{arccot} \left( -\frac{\sqrt{3}}{3} \right)$

$\mbox{arccot} \left( -\frac{\sqrt{3}}{3} \right) = \answer{\frac{2\pi}{3}}$
\end{problem}

\begin{problem}
$\mbox{arccot} \left( 0 \right)$

\begin{solution}
$\mbox{arccot} \left( 0 \right) = \frac{\pi}{2}$
\end{solution}
\end{problem}

\begin{problem}
$\mbox{arccot} \left( \frac{\sqrt{3}}{3} \right)$

$\mbox{arccot} \left( \frac{\sqrt{3}}{3} \right) = \answer{\frac{\pi}{3}}$
\end{problem}

\begin{problem}
$\mbox{arccot} \left( 1 \right)$

\begin{solution}
$\mbox{arccot} \left( 1 \right) = \frac{\pi}{4}$
\end{solution}
\end{problem}

\begin{problem}
$\mbox{arccot} \left( \sqrt{3} \right)$

$\mbox{arccot} \left( \sqrt{3} \right) = \answer{\frac{\pi}{6}}$
\end{problem}

\begin{problem}
$\mbox{arcsec} \left( 2 \right)$

\begin{solution}
$\mbox{arcsec} \left( 2 \right) = \frac{\pi}{3}$
\end{solution}
\end{problem}

\begin{problem}
$\mbox{arccsc} \left( 2 \right)$

$\mbox{arccsc} \left( 2 \right) = \answer{\frac{\pi}{6}}$
\end{problem}

\begin{problem}
$\mbox{arcsec} \left( \sqrt{2} \right)$

\begin{solution}
$\mbox{arcsec} \left( \sqrt{2} \right) = \frac{\pi}{4}$
\end{solution}
\end{problem}

\begin{problem}
$\mbox{arccsc} \left( \sqrt{2} \right)$

$\mbox{arccsc} \left( \sqrt{2} \right) = \answer{\frac{\pi}{4}}$
\end{problem}


\begin{problem}
$\mbox{arcsec} \left( \frac{2\sqrt{3}}{3} \right)$

\begin{solution}
$\mbox{arcsec} \left( \frac{2\sqrt{3}}{3} \right) = \frac{\pi}{6}$
\end{solution}
\end{problem}

\begin{problem}
$\mbox{arccsc} \left( \frac{2\sqrt{3}}{3} \right)$

$\mbox{arccsc} \left( \frac{2\sqrt{3}}{3} \right) = \answer{\frac{\pi}{3}}$
\end{problem}

\begin{problem}
$\mbox{arcsec} \left( 1 \right)$

\begin{solution}
$\mbox{arcsec} \left( 1 \right) = 0$
\end{solution}
\end{problem}

\begin{problem} \label{exactvaluearclast}
$\mbox{arccsc} \left( 1 \right)$

$\mbox{arccsc} \left( 1 \right) = \answer{\frac{\pi}{2}}$
\end{problem} 

\end{question}

\begin{question}

In Exercises \ref{trigfriendexactfirst} - \ref{trigfriendexactlast}, assume that the range of arcsecant is $\left[0, \frac{\pi}{2} \right) \cup \left( \frac{\pi}{2}, \pi \right]$ and that the range of arccosecant is
$\left[ -\frac{\pi}{2}, 0 \right)  \cup \left(0, \frac{\pi}{2} \right]$ when finding the exact value. (See Section \ref{arcsecanttrigfriendly}.)

\begin{problem} \label{trigfriendexactfirst}
$\mbox{arcsec} \left( -2 \right)$

\begin{solution}
$\mbox{arcsec} \left( -2 \right) = \frac{2\pi}{3}$
\end{solution}
\end{problem}

\begin{problem}
$\mbox{arcsec} \left( -\sqrt{2} \right)$ 

$\mbox{arcsec} \left( -\sqrt{2} \right) = \answer{\frac{3\pi}{4}}$
\end{problem}

\begin{problem}
$\mbox{arcsec} \left( -\frac{2\sqrt{3}}{3} \right)$

\begin{solution}
$\mbox{arcsec} \left( -\frac{2\sqrt{3}}{3} \right) = \frac{5\pi}{6}$
\end{solution}
\end{problem}

\begin{problem}
$\mbox{arcsec} \left( -1 \right)$

$\mbox{arcsec} \left( -1 \right) = \answer{\pi}$
\end{problem}

\begin{problem}
$\mbox{arccsc} \left( -2 \right)$

\begin{solution}
$\mbox{arccsc} \left( -2 \right) = -\frac{\pi}{6}$
\end{solution}
\end{problem}

\begin{problem}
$\mbox{arccsc} \left( -\sqrt{2} \right)$ 

$\mbox{arccsc} \left( -\sqrt{2} \right) = \answer{-\frac{\pi}{4}}$
\end{problem}

\begin{problem}
$\mbox{arccsc} \left( -\frac{2\sqrt{3}}{3} \right)$

\begin{solution}
$\mbox{arccsc} \left( -\frac{2\sqrt{3}}{3} \right) = -\frac{\pi}{3}$
\end{solution}
\end{problem}

\begin{problem} \label{trigfriendexactlast}
$\mbox{arccsc} \left( -1 \right)$

$\mbox{arccsc} \left( -1 \right) = \answer{-\frac{\pi}{2}}$
\end{problem}

\end{question}

\begin{question}
In Exercises \ref{calcfriendexactfirst} - \ref{calcfriendexactlast}, assume that the range of arcsecant is $\left[0, \frac{\pi}{2} \right) \cup \left[\pi, \frac{3\pi}{2} \right)$ and that the range of arccosecant is $\left(0, \frac{\pi}{2} \right] \cup \left( \pi, \frac{3\pi}{2} \right]$ when finding the exact value.  (See Section \ref{arcsecantcalcfriendly}.)


\begin{problem} \label{calcfriendexactfirst}
$\mbox{arcsec} \left( -2 \right)$

\begin{solution}
$\mbox{arcsec} \left( -2 \right) = \frac{4\pi}{3}$
\end{solution}
\end{problem}

\begin{problem}
$\mbox{arcsec} \left( -\sqrt{2} \right)$

$\mbox{arcsec} \left( -\sqrt{2} \right) = \answer{\frac{5\pi}{4}}$
\end{problem}

\begin{problem}
$\mbox{arcsec} \left( -\frac{2\sqrt{3}}{3} \right)$

\begin{solution}
$\mbox{arcsec} \left( -\frac{2\sqrt{3}}{3} \right) = \frac{7\pi}{6}$
\end{solution}
\end{problem}

\begin{problem}
$\mbox{arcsec} \left( -1 \right)$

$\mbox{arcsec} \left( -1 \right) = \answer{\pi}$
\end{problem}

\begin{problem}
$\mbox{arccsc} \left( -2 \right)$

\begin{solution}
$\mbox{arccsc} \left( -2 \right) = \frac{7\pi}{6}$
\end{solution}
\end{problem}

\begin{problem}
$\mbox{arccsc} \left( -\sqrt{2} \right)$ 

$\mbox{arccsc} \left( -\sqrt{2} \right) = \answer{\frac{5\pi}{4}}$
\end{problem}

\begin{problem}
$\mbox{arccsc} \left( -\frac{2\sqrt{3}}{3} \right)$

\begin{solution}
$\mbox{arccsc} \left( -\frac{2\sqrt{3}}{3} \right) = \frac{4\pi}{3}$
\end{solution}
\end{problem}

\begin{problem} \label{calcfriendexactlast}
$\mbox{arccsc} \left( -1 \right)$

$\mbox{arccsc} \left( -1 \right) = \answer{\frac{3\pi}{2}}$
\end{problem}

\end{question}

\begin{question}
In Exercises \ref{comboexactfirst} - \ref{comboexactlast}, find the exact value or state that it is undefined.

\begin{problem} \label{comboexactfirst}
$\sin\left(\arcsin\left(\frac{1}{2}\right)\right)$

\begin{solution}
$\sin\left(\arcsin\left(\frac{1}{2}\right)\right) = \frac{1}{2}$
\end{solution}
\end{problem}

\begin{problem}
$\sin\left(\arcsin\left(-\frac{\sqrt{2}}{2}\right)\right)$

$\sin\left(\arcsin\left(-\frac{\sqrt{2}}{2}\right)\right) = \answer{-\frac{\sqrt{2}}{2}}$
\end{problem}

\begin{problem}
$\sin\left(\arcsin\left(\frac{3}{5}\right)\right)$

\begin{solution}
$\sin\left(\arcsin\left(\frac{3}{5}\right)\right) = \frac{3}{5}$
\end{solution}
\end{problem}

\begin{problem}
$\sin\left(\arcsin\left(-0.42\right)\right)$

$\sin\left(\arcsin\left(-0.42\right)\right) = \answer{-0.42}$
\end{problem}

\begin{problem}
$\sin\left(\arcsin\left(\frac{5}{4}\right)\right)$

\begin{solution}
$\sin\left(\arcsin\left(\frac{5}{4}\right)\right)$ is undefined.
\end{solution}
\end{problem}

\begin{problem}
$\cos\left(\arccos\left(\frac{\sqrt{2}}{2}\right)\right)$

$\cos\left(\arccos\left(\frac{\sqrt{2}}{2}\right)\right) = \answer{\frac{\sqrt{2}}{2}}$
\end{problem}

\begin{problem}
$\cos\left(\arccos\left(-\frac{1}{2}\right)\right)$

\begin{solution}
$\cos\left(\arccos\left(-\frac{1}{2}\right)\right) = -\frac{1}{2}$
\end{solution}
\end{problem}

\begin{problem}
$\cos\left(\arccos\left(\frac{5}{13}\right)\right)$

$\cos\left(\arccos\left(\frac{5}{13}\right)\right) = \answer{\frac{5}{13}}$
\end{problem}

\begin{problem}
$\cos\left(\arccos\left(-0.998\right)\right)$

\begin{solution}
$\cos\left(\arccos\left(-0.998\right)\right) = -0.998$
\end{solution}
\end{problem}

\begin{problem}
$\cos\left(\arccos\left(\pi \right)\right)$

\begin{solution}
$\cos\left(\arccos\left(\pi \right)\right)$ is undefined.
\end{solution}
\end{problem}

\begin{problem}
$\tan\left(\arctan\left(-1\right)\right)$

\begin{solution}
$\tan\left(\arctan\left(-1\right)\right) = -1$
\end{solution}
\end{problem}

\begin{problem}
$\tan\left(\arctan\left(\sqrt{3}\right)\right)$

$\tan\left(\arctan\left(\sqrt{3}\right)\right) = \answer{\sqrt{3}}$
\end{problem}

\begin{problem}
 $\tan\left(\arctan\left(\frac{5}{12}\right)\right)$

\begin{solution}
$\tan\left(\arctan\left(\frac{5}{12}\right)\right) = \frac{5}{12}$
\end{solution}
\end{problem}

\begin{problem}
$\tan\left(\arctan\left(0.965\right)\right)$

$\tan\left(\arctan\left(0.965\right)\right) = \answer{0.965}$
\end{problem}

\begin{problem}
$\tan\left(\arctan\left( 3\pi \right)\right)$

\begin{solution}
$\tan\left(\arctan\left( 3\pi \right)\right) = 3\pi$
\end{solution}
\end{problem}

\begin{problem}
$\cot\left(\text{arccot}\left(1\right)\right)$

$\cot\left(\text{arccot}\left(1\right)\right) = \answer{1}$
\end{problem}

\begin{problem}
$\cot\left(\text{arccot}\left(-\sqrt{3}\right)\right)$

\begin{solution}
$\cot\left(\text{arccot}\left(-\sqrt{3}\right)\right) = -\sqrt{3}$ 
\end{solution}
\end{problem}

\begin{problem}
$\cot\left(\text{arccot}\left(-\frac{7}{24}\right)\right)$

$\cot\left(\text{arccot}\left(-\frac{7}{24}\right)\right) = \answer{-\frac{7}{24}}$
\end{problem}

\begin{problem}
$\cot\left(\text{arccot}\left(-0.001\right)\right)$

\begin{solution}
$\cot\left(\text{arccot}\left(-0.001\right)\right) = -0.001$
\end{solution}
\end{problem}

\begin{problem}
$\cot\left(\text{arccot}\left( \frac{17\pi}{4} \right)\right)$

$\cot\left(\text{arccot}\left( \frac{17\pi}{4} \right)\right) = \answer{\frac{17\pi}{4}}$
\end{problem}

\begin{problem}
$\sec\left(\text{arcsec}\left(2\right)\right)$

\begin{solution}
$\sec\left(\text{arcsec}\left(2\right)\right) = 2$
\end{solution}
\end{problem}

\begin{problem}
$\sec\left(\text{arcsec}\left(-1\right)\right)$

$\sec\left(\text{arcsec}\left(-1\right)\right) = \answer{-1}$
\end{problem}

\begin{problem}
$\sec\left(\text{arcsec}\left(\frac{1}{2}\right)\right)$

\begin{solution}
$\sec\left(\text{arcsec}\left(\frac{1}{2}\right)\right)$ is undefined.
\end{solution}
\end{problem}

\begin{problem}
$\sec\left(\text{arcsec}\left(0.75\right)\right)$

\begin{solution}
$\sec\left(\text{arcsec}\left(0.75\right)\right)$ is undefined.
\end{solution}
\end{problem}

\begin{problem}
$\sec\left(\text{arcsec}\left( 117\pi \right)\right)$

\begin{solution}
$\sec\left(\text{arcsec}\left( 117\pi \right)\right)= 117\pi $
\end{solution}
\end{problem}

\begin{problem}
$\csc\left(\text{arccsc}\left(\sqrt{2}\right)\right)$

$\csc\left(\text{arccsc}\left(\sqrt{2}\right)\right) = \answer{\sqrt{2}}$ 
\end{problem}

\begin{problem}
$\csc\left(\text{arccsc}\left(-\frac{2\sqrt{3}}{3}\right)\right)$

\begin{solution}
$\csc\left(\text{arccsc}\left(-\frac{2\sqrt{3}}{3}\right)\right) = -\frac{2\sqrt{3}}{3}$
\end{solution}
\end{problem}

\begin{problem}
$\csc\left(\text{arccsc}\left(\frac{\sqrt{2}}{2}\right)\right)$

\begin{solution}
$\csc\left(\text{arccsc}\left(\frac{\sqrt{2}}{2}\right)\right)$ is undefined.
\end{solution}
\end{problem}

\begin{problem}
$\csc\left(\text{arccsc}\left(1.0001\right)\right)$

\begin{solution}
$\csc\left(\text{arccsc}\left(1.0001\right)\right) = 1.0001$
\end{solution}
\end{problem}

\begin{problem} \label{comboexactlast}
$\csc\left(\text{arccsc}\left( \frac{\pi}{4} \right)\right)$

\begin{solution}
$\csc\left(\text{arccsc}\left( \frac{\pi}{4} \right)\right)$ is undefined.
\end{solution}
\end{problem}

\end{question}

\begin{question}
In Exercises \ref{morecomboexactfirst} - \ref{morecomboexactlast}, find the exact value or state that it is undefined.

\begin{problem} \label{morecomboexactfirst}
$\arcsin\left(\sin\left(\frac{\pi}{6}\right) \right)$ 

\begin{solution}
$\arcsin\left(\sin\left(\frac{\pi}{6}\right) \right) = \frac{\pi}{6}$
\end{solution}
\end{problem}

\begin{problem}
$\arcsin\left(\sin\left(-\frac{\pi}{3}\right) \right)$

$\arcsin\left(\sin\left(-\frac{\pi}{3}\right) \right) = \answer{-\frac{\pi}{3}}$
\end{problem}

\begin{problem}
$\arcsin\left(\sin\left(\frac{3\pi}{4}\right) \right)$

\begin{solution}
$\arcsin\left(\sin\left(\frac{3\pi}{4}\right) \right) = \frac{\pi}{4}$
\end{solution}
\end{problem}

\begin{problem}
$\arcsin\left(\sin\left(\frac{11\pi}{6}\right) \right)$

$\arcsin\left(\sin\left(\frac{11\pi}{6}\right) \right) = \answer{-\frac{\pi}{6}}$
\end{problem}

\begin{problem}
$\arcsin\left(\sin\left(\frac{4\pi}{3}\right) \right)$

\begin{solution}
$\arcsin\left(\sin\left(\frac{4\pi}{3}\right) \right) = -\frac{\pi}{3}$
\end{solution}
\end{problem}

\begin{problem}
$\arccos\left(\cos\left(\frac{\pi}{4}\right) \right)$

$\arccos\left(\cos\left(\frac{\pi}{4}\right) \right) = \answer{\frac{\pi}{4}}$
\end{problem}

\begin{problem} 
$\arccos\left(\cos\left(\frac{2\pi}{3}\right) \right)$

\begin{solution}
$\arccos\left(\cos\left(\frac{2\pi}{3}\right) \right) = \frac{2\pi}{3}$
\end{solution}
\end{problem}

\begin{problem}
$\arccos\left(\cos\left(\frac{3\pi}{2}\right) \right)$

$\arccos\left(\cos\left(\frac{3\pi}{2}\right) \right) = \answer{\frac{\pi}{2}}$
\end{problem}

\begin{problem}
$\arccos\left(\cos\left(-\frac{\pi}{6}\right) \right)$

\begin{solution}
$\arccos\left(\cos\left(-\frac{\pi}{6}\right) \right) = \frac{\pi}{6}$ 
\end{solution}
\end{problem}

\begin{problem}
$\arccos\left(\cos\left(\frac{5\pi}{4}\right) \right)$

$\arccos\left(\cos\left(\frac{5\pi}{4}\right) \right) = \answer{\frac{3\pi}{4}}$
\end{problem}

\begin{problem}
$\arctan\left(\tan\left(\frac{\pi}{3}\right) \right)$

\begin{solution}
$\arctan\left(\tan\left(\frac{\pi}{3}\right) \right) = \frac{\pi}{3}$
\end{solution}
\end{problem}

\begin{problem}
$\arctan\left(\tan\left(-\frac{\pi}{4}\right) \right)$ 

$\arctan\left(\tan\left(-\frac{\pi}{4}\right) \right) = \answer{-\frac{\pi}{4}}$
\end{problem}

\begin{problem}
$\arctan\left(\tan\left(\pi\right) \right)$

\begin{solution}
$\arctan\left(\tan\left(\pi\right) \right) = 0$
\end{solution}
\end{problem}

\begin{problem}
$\arctan\left(\tan\left(\frac{\pi}{2}\right) \right)$

\begin{solution}
$\arctan\left(\tan\left(\frac{\pi}{2}\right) \right)$ is undefined
\end{solution}
\end{problem}

\begin{problem}
$\arctan\left(\tan\left(\frac{2\pi}{3}\right) \right)$

\begin{solution}
$\arctan\left(\tan\left(\frac{2\pi}{3}\right) \right) = -\frac{\pi}{3}$
\end{solution}
\end{problem}

\begin{problem}  $\text{arccot}\left(\cot\left(\frac{\pi}{3}\right) \right)$

$\text{arccot}\left(\cot\left(\frac{\pi}{3}\right) \right) = \answer{\frac{\pi}{3}}$
\end{problem}

\begin{problem}
$\text{arccot}\left(\cot\left(-\frac{\pi}{4}\right) \right)$

\begin{solution}
$\text{arccot}\left(\cot\left(-\frac{\pi}{4}\right) \right) = \frac{3\pi}{4}$
\end{solution}
\end{problem}

\begin{problem}
$\text{arccot}\left(\cot\left(\pi\right) \right)$

\begin{solution}
$\text{arccot}\left(\cot\left(\pi\right) \right)$ is undefined
\end{solution}
\end{problem}

\begin{problem}
$\text{arccot}\left(\cot\left(\frac{3\pi}{2}\right) \right)$

\begin{solution}
$\text{arccot}\left(\cot\left(\frac{3\pi}{2}\right) \right) = \frac{\pi}{2}$
\end{solution}
\end{problem}

\begin{problem} \label{morecomboexactlast}
$\text{arccot}\left(\cot\left(\frac{2\pi}{3}\right) \right)$ 

$\text{arccot}\left(\cot\left(\frac{2\pi}{3}\right) \right) = \answer{\frac{2\pi}{3}}$
\end{problem}

\end{question}

\begin{question}

In Exercises \ref{moreextracombofirst} - \ref{moreextracombolast}, assume that the range of arcsecant is $\left[0, \frac{\pi}{2} \right) \cup \left( \frac{\pi}{2}, \pi \right]$ and that the range of arccosecant is $\left[ -\frac{\pi}{2}, 0 \right)  \cup \left(0, \frac{\pi}{2} \right]$ when finding the exact value. (See Section \ref{arcsecanttrigfriendly}.)

\begin{problem} \label{moreextracombofirst}
$\text{arcsec}\left(\sec\left(\frac{\pi}{4}\right) \right)$

\begin{solution}
$\text{arcsec}\left(\sec\left(\frac{\pi}{4}\right) \right) = \frac{\pi}{4}$ 
\end{solution}
\end{problem}

\begin{problem}
$\text{arcsec}\left(\sec\left(\frac{4\pi}{3}\right) \right)$

$\text{arcsec}\left(\sec\left(\frac{4\pi}{3}\right) \right) = \answer{\frac{2\pi}{3}}$
\end{problem}

\begin{problem}
$\text{arcsec}\left(\sec\left( \frac{5\pi}{6} \right) \right)$

\begin{solution}
$\text{arcsec}\left(\sec\left( \frac{5\pi}{6} \right) \right) = \frac{5\pi}{6}$
\end{solution}
\end{problem}

 \begin{problem}
 $\text{arcsec}\left(\sec\left(-\frac{\pi}{2} \right) \right)$

 \begin{solution}
 $\text{arcsec}\left(\sec\left(-\frac{\pi}{2} \right) \right)$ is undefined.
 \end{solution}
 \end{problem}
 
\begin{problem}
$\text{arcsec}\left(\sec\left(\frac{5\pi}{3}\right) \right)$

\begin{solution}
$\text{arcsec}\left(\sec\left(\frac{5\pi}{3}\right) \right) = \frac{\pi}{3}$
\end{solution}
\end{problem}

\begin{problem}
$\text{arccsc}\left(\csc\left(\frac{\pi}{6}\right) \right)$

$\text{arccsc}\left(\csc\left(\frac{\pi}{6}\right) \right) = \answer{\frac{\pi}{6}}$
\end{problem}

\begin{problem}
$\text{arccsc}\left(\csc\left(\frac{5\pi}{4}\right) \right)$

\begin{solution}
$\text{arccsc}\left(\csc\left(\frac{5\pi}{4}\right) \right) = -\frac{\pi}{4}$
\end{solution}
\end{problem}

\begin{problem}
$\text{arccsc}\left(\csc\left( \frac{2\pi}{3} \right) \right)$

$\text{arccsc}\left(\csc\left( \frac{2\pi}{3} \right) \right) = \answer{\frac{\pi}{3}}$
\end{problem}

\begin{problem}
$\text{arccsc}\left(\csc\left(-\frac{\pi}{2} \right) \right)$

\begin{solution}
$\text{arccsc}\left(\csc\left(-\frac{\pi}{2} \right) \right) = -\frac{\pi}{2}$
\end{solution}
\end{problem}

\begin{problem}
$\text{arccsc}\left(\csc\left(\frac{11\pi}{6}\right) \right)$

$\text{arccsc}\left(\csc\left(\frac{11\pi}{6}\right) \right) = \answer{-\frac{\pi}{6}}$
\end{problem}

\begin{problem}
$\text{arcsec}\left(\sec\left(\frac{11\pi}{12}\right) \right)$

\begin{solution}
$\text{arcsec}\left(\sec\left(\frac{11\pi}{12}\right) \right) = \frac{11\pi}{12}$
\end{solution}
\end{problem}

\begin{problem} \label{moreextracombolast}
$\text{arccsc}\left(\csc\left(\frac{9\pi}{8}\right) \right)$ 

$\text{arccsc}\left(\csc\left(\frac{9\pi}{8}\right) \right) = \answer{-\frac{\pi}{8}}$
\end{problem}

\end{question}

\begin{question}
In Exercises \ref{extracombofirst} - \ref{extracombolast}, assume that the range of arcsecant is $\left[0, \frac{\pi}{2} \right) \cup \left[\pi, \frac{3\pi}{2} \right)$ and that the range of arccosecant is $\left(0, \frac{\pi}{2} \right] \cup \left( \pi, \frac{3\pi}{2} \right]$ when finding the exact value. (See Section \ref{arcsecantcalcfriendly}.)


\begin{problem} \label{extracombofirst}
$\text{arcsec}\left(\sec\left(\frac{\pi}{4}\right) \right)$

\begin{solution}
$\text{arcsec}\left(\sec\left(\frac{\pi}{4}\right) \right) = \frac{\pi}{4}$
\end{solution}
\end{problem}

\begin{problem}
$\text{arcsec}\left(\sec\left(\frac{4\pi}{3}\right) \right)$

$\text{arcsec}\left(\sec\left(\frac{4\pi}{3}\right) \right) = \answer{\frac{4\pi}{3}}$
\end{problem}

\begin{problem}
$\text{arcsec}\left(\sec\left( \frac{5\pi}{6} \right) \right)$

\begin{solution}
$\text{arcsec}\left(\sec\left( \frac{5\pi}{6} \right) \right) = \frac{7\pi}{6}$
\end{solution}
\end{problem}

\begin{problem}
$\text{arcsec}\left(\sec\left(-\frac{\pi}{2} \right) \right)$

\begin{solution}
$\text{arcsec}\left(\sec\left(-\frac{\pi}{2} \right) \right)$ is undefined.
\end{solution}
\end{problem}

\begin{problem}
$\text{arcsec}\left(\sec\left(\frac{5\pi}{3}\right) \right)$

\begin{solution}
$\text{arcsec}\left(\sec\left(\frac{5\pi}{3}\right) \right) = \frac{\pi}{3}$
\end{solution}
\end{problem}

\begin{problem}
$\text{arccsc}\left(\csc\left(\frac{\pi}{6}\right) \right)$

$\text{arccsc}\left(\csc\left(\frac{\pi}{6}\right) \right) = \answer{\frac{\pi}{6}}$
\end{problem}

\begin{problem}  $\text{arccsc}\left(\csc\left(\frac{5\pi}{4}\right) \right)$

\begin{solution}
$\text{arccsc}\left(\csc\left(\frac{5\pi}{4}\right) \right) = \frac{5\pi}{4}$
\end{solution}
\end{problem}

\begin{problem}
$\text{arccsc}\left(\csc\left( \frac{2\pi}{3} \right) \right)$

$\text{arccsc}\left(\csc\left( \frac{2\pi}{3} \right) \right) = \answer{\frac{\pi}{3}}$
\end{problem}

\begin{problem}
$\text{arccsc}\left(\csc\left(-\frac{\pi}{2} \right) \right)$

\begin{solution}
$\text{arccsc}\left(\csc\left(-\frac{\pi}{2} \right) \right) = \frac{3\pi}{2}$
\end{solution}
\end{problem}

\setcounter{HW}{\value{enumi}}

\begin{problem}
$\text{arccsc}\left(\csc\left(\frac{11\pi}{6}\right) \right)$

$\text{arccsc}\left(\csc\left(\frac{11\pi}{6}\right) \right) = \answer{\frac{7\pi}{6}}$
\end{problem}

\begin{problem}
$\text{arcsec}\left(\sec\left(\frac{11\pi}{12}\right) \right)$

\begin{solution}
$\text{arcsec}\left(\sec\left(\frac{11\pi}{12}\right) \right) = \frac{13\pi}{12}$
\end{solution}
\end{problem}

\begin{problem} \label{extracombolast}
$\text{arccsc}\left(\csc\left(\frac{9\pi}{8}\right) \right)$ 

$\text{arccsc}\left(\csc\left(\frac{9\pi}{8}\right) \right) = \answer{\frac{9\pi}{8}}$
\end{problem}

\end{question}

\begin{question}
In Exercises \ref{stillmoreexactfirst} - \ref{stillmoreexactlast}, find the exact value or state that it is undefined.

\begin{problem}  \label{stillmoreexactfirst}
$\sin\left(\arccos\left(-\frac{1}{2}\right)\right)$

\begin{solution}
$\sin\left(\arccos\left(-\frac{1}{2}\right)\right) = \frac{\sqrt{3}}{2}$
\end{solution}
\end{problem}

\begin{problem}
$\sin\left(\arccos\left(\frac{3}{5}\right)\right)$

$\sin\left(\arccos\left(\frac{3}{5}\right)\right) = \answer{\frac{4}{5}}$
\end{problem}

\begin{problem}
$\sin\left(\arctan\left(-2\right)\right)$ 

\begin{solution}
$\sin\left(\arctan\left(-2\right)\right) = -\frac{2\sqrt{5}}{5}$
\end{solution}
\end{problem}

\begin{problem}  $\sin\left(\text{arccot}\left(\sqrt{5}\right)\right)$

$\sin\left(\text{arccot}\left(\sqrt{5}\right)\right) = \answer{\frac{\sqrt{6}}{6}}$
\end{problem}

\begin{problem}
$\sin\left(\text{arccsc}\left(-3\right)\right)$

\begin{solution}
$\sin\left(\text{arccsc}\left(-3\right)\right) = -\frac{1}{3}$
\end{solution}
\end{problem}

\begin{problem}
$\cos\left(\arcsin\left(-\frac{5}{13}\right)\right)$

$\cos\left(\arcsin\left(-\frac{5}{13}\right)\right) = \answer{\frac{12}{13}}$
\end{problem}

\begin{problem}
$\cos\left(\arctan\left(\sqrt{7} \right)\right)$

\begin{solution}
$\cos\left(\arctan\left(\sqrt{7} \right)\right) = \frac{\sqrt{2}}{4}$
\end{solution}
\end{problem}

\begin{problem}
$\cos\left(\text{arccot}\left( 3 \right)\right)$

$\cos\left(\text{arccot}\left( 3 \right)\right) = \answer{\frac{3\sqrt{10}}{10}}$
\end{problem}

\begin{problem}
$\cos\left(\text{arcsec}\left( 5 \right)\right)$

\begin{solution}
$\cos\left(\text{arcsec}\left( 5 \right)\right) = \frac{1}{5}$ 
\end{solution}
\end{problem}

\begin{problem}
$\tan\left(\arcsin\left(-\frac{2\sqrt{5}}{5}\right)\right)$

$\tan\left(\arcsin\left(-\frac{2\sqrt{5}}{5}\right)\right)=\answer{-2}$
\end{problem}

\begin{problem}
$\tan\left(\arccos\left(-\frac{1}{2}\right)\right)$ 

\begin{solution}
$\tan\left(\arccos\left(-\frac{1}{2}\right)\right) = -\sqrt{3}$
\end{solution}
\end{problem}

\begin{problem}
$\tan\left(\text{arcsec}\left(\frac{5}{3}\right)\right)$

$\tan\left(\text{arcsec}\left(\frac{5}{3}\right)\right) = \answer{\frac{4}{3}}$
\end{problem}

\begin{problem}
$\tan\left(\text{arccot}\left( 12  \right)\right)$

\begin{solution}
$\tan\left(\text{arccot}\left( 12  \right)\right) = \frac{1}{12}$
\end{solution}
\end{problem}

\begin{problem}
$\cot\left(\arcsin\left(\frac{12}{13}\right)\right)$

$\cot\left(\arcsin\left(\frac{12}{13}\right)\right) = \answer{\frac{5}{12}}$
\end{problem}

\begin{problem}
$\cot\left(\arccos\left(\frac{\sqrt{3}}{2}\right)\right)$

\begin{solution}
$\cot\left(\arccos\left(\frac{\sqrt{3}}{2}\right)\right) = \sqrt{3}$
\end{solution}
\end{problem}

\begin{problem}  $\cot\left(\text{arccsc}\left(\sqrt{5}\right)\right)$

$\cot\left(\text{arccsc}\left(\sqrt{5}\right)\right) = \answer{2}$
\end{problem}

\begin{problem}
$\cot\left(\arctan \left( 0.25 \right)\right)$

\begin{solution}
$\cot\left(\arctan \left( 0.25 \right)\right) = 4$
\end{solution}
\end{problem}

\begin{problem}
$\sec\left(\arccos\left(\frac{\sqrt{3}}{2}\right)\right)$

$\sec\left(\arccos\left(\frac{\sqrt{3}}{2}\right)\right) = \answer{\frac{2\sqrt{3}}{3}}$
\end{problem}

\begin{problem}
$\sec\left(\arcsin\left(-\frac{12}{13}\right)\right)$ 

\begin{solution}
$\sec\left(\arcsin\left(-\frac{12}{13}\right)\right) = \frac{13}{5}$
\end{solution}
\end{problem}

\begin{problem}
$\sec\left(\arctan\left(10\right)\right)$

$\sec\left(\arctan\left(10\right)\right) = \answer{\sqrt{101}}$ 
\end{problem}

\begin{problem} $\sec\left(\text{arccot}\left(-\frac{\sqrt{10}}{10}\right)\right)$

\begin{solution}
$\sec\left(\text{arccot}\left(-\frac{\sqrt{10}}{10}\right)\right) = -\sqrt{11}$
\end{solution}
\end{problem}

\begin{problem}
$\csc\left(\text{arccot}\left(9 \right)\right)$

$\csc\left(\text{arccot}\left(9 \right)\right) = \answer{\sqrt{82}}$
\end{problem}

\begin{problem}
$\csc\left(\arcsin\left(\frac{3}{5}\right)\right)$

\begin{solution}
$\csc\left(\arcsin\left(\frac{3}{5}\right)\right) = \frac{5}{3}$
\end{solution}
\end{problem}

\begin{problem} \label{stillmoreexactlast}
$\csc\left(\arctan\left(-\frac{2}{3}\right)\right)$

$\csc\left(\arctan\left(-\frac{2}{3}\right)\right) = \answer{-\frac{\sqrt{13}}{2}}$
\end{problem}

\end{question}

\begin{question}
In Exercises \ref{exactvalueidenfirst} - \ref{exactvalueidenlast}, find the exact value or state that it is undefined.

\begin{problem} \label{exactvalueidenfirst}
$\sin\left(\arcsin\left( \frac{5}{13} \right) + \frac{\pi}{4}\right)$

\begin{solution}
$\sin\left(\arcsin\left( \frac{5}{13} \right) + \frac{\pi}{4}\right) = \frac{17\sqrt{2}}{26}$
\end{solution}
\end{problem}

\begin{problem}
$\cos\left( \text{arcsec}(3) + \arctan(2) \right)$

$\cos\left( \text{arcsec}(3) + \arctan(2) \right) = \answer{\frac{\sqrt{5} - 4\sqrt{10}}{15}}$
\end{problem}

\begin{problem}
$\tan\left( \arctan(3) + \arccos\left(-\frac{3}{5}\right) \right)$

\begin{solution}
$\tan\left( \arctan(3) + \arccos\left(-\frac{3}{5}\right) \right) = \frac{1}{3}$
\end{solution}
\end{problem}

\begin{problem}
$\sin\left(2 \arcsin\left(-\frac{4}{5}\right)\right)$

$\sin\left(2 \arcsin\left(-\frac{4}{5}\right)\right)= \answer{-\frac{24}{25}}$
\end{problem}

\begin{problem}
$\sin\left(2\text{arccsc}\left(\frac{13}{5}\right)\right)$

\begin{solution}
$\sin\left(2\text{arccsc}\left(\frac{13}{5}\right)\right) = \frac{120}{169}$
\end{solution}
\end{problem}

\begin{problem}
$\sin\left(2 \arctan\left(2\right)\right)$ 

$\sin\left(2 \arctan\left(2\right)\right) = \frac{4}{5}$
\end{problem}

\begin{problem}
$\cos\left(2 \arcsin\left(\frac{3}{5}\right)\right)$

\begin{solution}
$\cos\left(2 \arcsin\left(\frac{3}{5}\right)\right) = \frac{7}{25}$
\end{solution}
\end{problem}

\begin{problem}
$\cos\left(2 \text{arcsec}\left(\frac{25}{7}\right)\right)$

$\cos\left(2 \text{arcsec}\left(\frac{25}{7}\right)\right) = \answer{-\frac{527}{625}}$
\end{problem}

\begin{problem}
$\cos\left(2 \text{arccot}\left(-\sqrt{5}\right)\right)$ 

\begin{solution}
$\cos\left(2 \text{arccot}\left(-\sqrt{5}\right)\right) = \frac{2}{3}$ 
\end{solution}
\end{problem}

\begin{problem} \label{exactvalueidenlast}
$\sin\left( \frac{\arctan(2)}{2} \right)$

$\sin\left( \frac{\arctan(2)}{2} \right) = \answer{\sqrt{\frac{5-\sqrt{5}}{10}}}$
\end{problem}

\end{question}

\begin{question}

In Exercises \ref{rewritefirst} - \ref{rewritelast}, rewrite each of the following composite functions as algebraic functions of $x$ and state the domain.

\begin{problem} \label{rewritefirst}
$f(x) = \sin \left( \arccos \left( x \right) \right)$ 

\begin{solution}
$f(x) = \sin \left( \arccos \left( x \right) \right) = \sqrt{1 - x^{2}}$  for $-1 \leq x \leq 1$
\end{solution}
\end{problem}

\begin{problem}
$f(x) = \cos \left( \arctan \left( x \right) \right)$

\begin{solution}
$f(x) = \cos \left( \arctan \left( x \right) \right) = \frac{1}{\sqrt{1 + x^{2}}}$ for all $x$
\end{solution}
\end{problem}

\begin{problem}
$f(x) = \tan \left( \arcsin \left( x \right) \right)$

\begin{solution}
$f(x) =\tan \left( \arcsin \left( x \right) \right) = \frac{x}{\sqrt{1 - x^{2}}}$  for $-1 < x < 1$
\end{solution}
\end{problem}

\begin{problem}
$f(x) = \sec \left( \arctan \left( x \right) \right)$

\begin{solution}
$f(x) =\sec \left( \arctan \left( x \right) \right) = \sqrt{1 + x^{2}}$ for all $x$
\end{solution}
\end{problem}

\begin{problem}
$f(x) = \csc \left( \arccos \left( x \right) \right)$

\begin{solution}
$f(x) =\csc \left( \arccos \left( x \right) \right) = \frac{1}{\sqrt{1 - x^{2}}}$ for $-1 < x < 1$
\end{solution}
\end{problem}

\begin{problem}
$f(x) = \sin \left( 2 \arctan \left( x \right) \right)$ 

\begin{solution}
$f(x) =\sin \left( 2 \arctan \left( x \right) \right) = \frac{2x}{x^{2} + 1}$ for all $x$
\end{solution}
\end{problem}

\begin{problem}
$f(x) = \sin \left( 2 \arccos \left( x \right) \right)$

$f(x) =\sin \left( 2 \arccos \left( x \right) \right) = 2x\sqrt{1-x^2}$  for $-1 \leq x \leq 1$
\end{problem}

\begin{problem}
$f(x) = \cos \left( 2 \arctan \left( x \right) \right)$ 

\begin{solution}
$f(x) =\cos \left( 2 \arctan \left( x \right) \right) = \frac{1 - x^{2}}{1 + x^{2}}$ for all $x$
\end{solution}
\end{problem}

\begin{problem}
$f(x) = \sin(\arccos(2x))$

\begin{solution}
$f(x) =\sin(\arccos(2x)) = \sqrt{1-4x^2}$ for $-\frac{1}{2} \leq x \leq \frac{1}{2}$
\end{solution}
\end{problem}

\begin{problem}
$f(x) = \sin\left(\arccos\left(\frac{x}{5}\right)\right)$

\begin{solution}
$f(x) =\sin\left(\arccos\left(\frac{x}{5}\right)\right) = \frac{\sqrt{25-x^2}}{5}$ for $-5 \leq x \leq 5$
\end{solution}
\end{problem}

\begin{problem}
$f(x) = \cos\left(\arcsin\left(\frac{x}{2}\right)\right)$

\begin{solution}
$f(x) =\cos\left(\arcsin\left(\frac{x}{2}\right)\right) = \frac{\sqrt{4-x^2}}{2}$ for $-2 \leq x \leq 2$
\end{solution}
\end{problem}

\begin{problem}
$f(x) = \cos\left(\arctan\left(3x\right)\right)$

\begin{solution}
$f(x) =\cos\left(\arctan\left(3x\right)\right) = \frac{1}{\sqrt{1+9x^{2}}}$ for all $x$
\end{solution}
\end{problem}

\begin{problem} $f(x) = \sin(2 \arcsin(7x))$

\begin{solution}
$f(x) =\sin(2 \arcsin(7x)) = 14x \sqrt{1-49x^2}$ for $-\frac{1}{7} \leq x \leq \frac{1}{7}$
\end{solution}
\end{problem}

\begin{problem}
$f(x) = \sin\left(2 \arcsin\left( \frac{x\sqrt{3}}{3} \right) \right)$

\begin{solution}
$f(x) =\sin\left(2 \arcsin\left( \frac{x\sqrt{3}}{3} \right) \right) = \frac{2x\sqrt{3-x^2}}{3}$ for $-\sqrt{3} \leq x \leq \sqrt{3}$
\end{solution}
\end{problem}

\begin{problem}
$f(x) = \cos(2 \arcsin(4x))$

\begin{solution}
$f(x) =\cos(2 \arcsin(4x)) = 1 - 32x^2$ for $-\frac{1}{4} \leq x \leq \frac{1}{4}$
\end{solution}
\end{problem}

\begin{problem}
$f(x) = \sec(\arctan(2x))\tan(\arctan(2x))$

\begin{solution}
$f(x) =\sec(\arctan(2x))\tan(\arctan(2x)) = 2x \sqrt{1+4x^2}$ for all $x$
\end{solution}
\end{problem}

\begin{problem}
$f(x) = \sin \left( \arcsin(x) + \arccos(x) \right)$

\begin{solution}
$f(x) =\sin \left( \arcsin(x) + \arccos(x) \right) = 1$ for $-1 \leq x \leq 1$
\end{solution}
\end{problem}

\begin{problem}
$f(x) = \cos \left( \arcsin(x) + \arctan(x) \right)$

\begin{solution}
$f(x) =\cos \left( \arcsin(x) + \arctan(x) \right) = \frac{\sqrt{1 - x^{2}} - x^{2}}{\sqrt{1 + x^{2}}}$ for $-1 \leq x \leq 1$
\end{solution}
\end{problem}

\begin{problem}
$f(x) = \tan \left( 2 \arcsin(x) \right)$ 

\begin{solution}
$f(x) =\tan \left( 2 \arcsin(x) \right) = \frac{2x\sqrt{1 - x^{2}}}{1 - 2x^{2}}$ for\footnote{The equivalence for $x = \pm 1$ can be verified independently of the derivation of the formula, but Calculus is required to fully understand what is happening at those $x$ values.  You'll see what we mean when you work through the details of the identity for $\tan(2t).$  For now, we exclude $x = \pm 1$ from our answer.}  $x$ in $\left(-1, -\frac{\sqrt{2}}{2}\right) \cup \left(-\frac{\sqrt{2}}{2}, \frac{\sqrt{2}}{2} \right) \cup \left(\frac{\sqrt{2}}{2}, 1\right)$
\end{solution}
\end{problem}

\begin{problem} \label{rewritelast}
$f(x) = \sin \left( \frac{1}{2}\arctan(x) \right)$ 

\begin{solution}
$f(x) =\sin \left( \frac{1}{2}\arctan(x) \right) = \left\{ \begin{array}{rr} \sqrt{\frac{\sqrt{x^{2} + 1} - 1}{2\sqrt{x^{2} + 1}}} & \text{for $x \geq 0$} \\ [10pt] -\sqrt{\frac{\sqrt{x^{2} + 1} - 1}{2\sqrt{x^{2} + 1}}} & \text{for $x < 0$}  \end{array}\right. $
\end{solution}
\end{problem}

\end{question}

\begin{problem}
If $\theta = \arcsin\left(\frac{x}{2}\right)$,  find an expression for $\theta + \sin(2\theta)$ in terms of $x$.

\begin{solution}
$\theta + \sin(2\theta) = \arcsin \left( \frac{x}{2} \right) + \frac{x\sqrt{4 - x^{2}}}{2}$
\end{solution}
\end{problem}

\begin{problem}
If $\theta = \arctan\left(\frac{x}{7}\right)$, find an expression for $\frac{1}{2}\theta - \frac{1}{2}\sin(2\theta)$ in terms of $x$.

\begin{solution}
$\frac{1}{2}\theta - \frac{1}{2}\sin(2\theta) = \frac{1}{2} \arctan \left( \frac{x}{7} \right) - \frac{7x}{x^{2} + 49}$
\end{solution}
\end{problem}


\begin{problem}
If $\theta = \mbox{arcsec}\left(\frac{x}{4}\right)$,  find an expression for $4\tan(\theta) - 4\theta$ in terms of $x$ assuming $x \geq 4$.

\begin{solution}
$4\tan(\theta) - 4\theta = \sqrt{x^{2} - 16} - 4\mbox{arcsec} \left( \frac{x}{4} \right)$
\end{solution}
\end{problem}


\begin{question}
In Exercises \ref{equarctrigfirst} - \ref{equarctriglast}, solve the equation using the techniques discussed in Example \ref{basicinverseeqns} then approximate the solutions which lie in the interval $[0, 2\pi)$ to four decimal places.

\begin{problem} \label{equarctrigfirst}
$\sin(\theta) = \frac{7}{11}$ 

\begin{solution}
$\theta = \arcsin\left(\frac{7}{11}\right) + 2\pi k$ or $\theta = \pi - \arcsin\left(\frac{7}{11}\right) + 2\pi k$, in  $[0, 2\pi)$, $\theta \approx 0.6898, \, 2.4518$
\end{solution}
\end{problem}

\begin{problem}
$\cos(\theta) = -\frac{2}{9}$

\begin{solution}
$\theta = \arccos\left(-\frac{2}{9}\right) + 2\pi k$ or $\theta = - \arccos\left(-\frac{2}{9}\right) + 2\pi k$, in  $[0, 2\pi)$, $\theta \approx 1.7949, \, 4.4883$
\end{solution}
\end{problem}

\begin{problem}
$\sin(\theta) = -0.569$

\begin{solution}
$\theta = \pi + \arcsin(0.569) + 2\pi k$ or $\theta = 2\pi - \arcsin(0.569) + 2\pi k$, in  $[0, 2\pi)$, $\theta \approx 3.7469, \, 5.6779$
\end{solution}
\end{problem}

\begin{problem}
$\cos(\theta) = 0.117$

\begin{solution}
$\theta= \arccos(0.117) + 2\pi k$ or $\theta = 2\pi - \arccos(0.117) + 2\pi k$, in  $[0, 2\pi)$, $\theta \approx 1.4535, \, 4.8297$
\end{solution}
\end{problem}

\begin{problem}
$\sin(\theta) = 0.008$

\begin{solution}
$\theta = \arcsin(0.008) + 2\pi k$ or $\theta = \pi - \arcsin(0.008) + 2\pi k$, in  $[0, 2\pi)$, $\theta \approx 0.0080, \, 3.1336$
\end{solution}
\end{problem}

\begin{problem}
$\cos(\theta) = \frac{359}{360}$

\begin{solution}
$\theta = \arccos\left(\frac{359}{360}\right) + 2\pi k$ or $\theta = 2\pi - \arccos\left(\frac{359}{360}\right) + 2\pi k$, in  $[0, 2\pi)$, $\theta \approx 0.0746, \, 6.2086$
\end{solution}
\end{problem}

\begin{problem}
$\tan(t) = 117$

\begin{solution}
$t = \arctan(117) + \pi k$, in  $[0, 2\pi)$, $t \approx 1.56225, \, 4.70384$
$t = \arctan\left(-\frac{1}{12}\right) + \pi k$, in  $[0, 2\pi)$,  $t \approx 3.0585, \, 6.2000$
\end{solution}
\end{problem}

\begin{problem}
$\cot(t) = -12$

\begin{solution}
$t = \arctan\left(-\frac{1}{12}\right) + \pi k$, in  $[0, 2\pi)$,  $t \approx 3.0585, \, 6.2000$
\end{solution}
\end{problem}

\begin{problem}
$\sec(t) = \frac{3}{2}$

\begin{solution}
$t = \arccos\left(\frac{2}{3}\right) + 2\pi k$ or $t = 2\pi - \arccos\left(\frac{2}{3}\right) + 2\pi k$, in  $[0, 2\pi)$, $t \approx 0.8411, \, 5.4422$
\end{solution}
\end{problem}

\begin{problem}
$\csc(t) = -\frac{90}{17}$

\begin{solution}
$t = \pi + \arcsin\left(\frac{17}{90}\right) + 2\pi k$ or $t = 2\pi - \arcsin\left(\frac{17}{90}\right) + 2\pi k$, in  $[0, 2\pi)$, $t \approx 3.3316, \, 6.0932$
\end{solution}
\end{problem}

\begin{problem}
$\tan(t) = -\sqrt{10}$

\begin{solution}
$t = \arctan\left(-\sqrt{10}\right) + \pi k$, in  $[0, 2\pi)$, $t \approx 1.8771, \, 5.0187$
\end{solution}
\end{problem}

\begin{problem}
$\sin(t) = \frac{3}{8}$

\begin{solution}
$t = \arcsin\left(\frac{3}{8}\right) + 2\pi k$ or $t = \pi - \arcsin\left(\frac{3}{8}\right) + 2\pi k$, in  $[0, 2\pi)$, $t \approx 0.3844, \, 2.7572$
\end{solution}
\end{problem}

\begin{problem}
$\cos(x) = -\frac{7}{16}$

\begin{solution}
$x =  \arccos\left(-\frac{7}{16}\right) + 2\pi k$ or $x = - \arccos\left(-\frac{7}{16}\right) + 2\pi k$, in  $[0, 2\pi)$, $x \approx 2.0236, \, 4.2596$
\end{solution}
\end{problem}

\begin{problem}
$\tan(x) = 0.03$

\begin{solution}
$x = \arctan(0.03) + \pi k$, in  $[0, 2\pi)$, $x \approx 0.0300, \, 3.1716$
\end{solution}
\end{problem}

\begin{problem}
$\sin(x) = 0.3502$

\begin{solution}
$x = \arcsin(0.3502) + 2\pi k$ or $x = \pi - \arcsin(0.3502) + 2\pi k$, in  $[0, 2\pi)$, $x \approx 0.3578, \,2.784$
\end{solution}
\end{problem}

\begin{problem}
$\sin(x) = -0.721$

\begin{solution}
$x = \pi + \arcsin(0.721) + 2\pi k$ or $x = 2\pi - \arcsin(0.721) + 2\pi k$, in  $[0, 2\pi)$, $x \approx 3.9468, \, 5.4780$
\end{solution}
\end{problem}

\begin{problem}
$\cos(x) = 0.9824$

\begin{solution}
$x = \arccos(0.9824) + 2\pi k$ or $x = 2\pi - \arccos(0.9824) + 2\pi k$, in  $[0, 2\pi)$, $x \approx 0.1879, \, 6.0953$
\end{solution}
\end{problem}

\begin{problem}
$\cos(x) = -0.5637$

\begin{solution}
$x = \arccos(-0.5637) + 2\pi k$ or $x = - \arccos(-0.5637)  + 2\pi k$, in  $[0, 2\pi)$, $x \approx 2.1697, \, 4.1135$
\end{solution}
\end{problem}

\begin{problem}
$\cot(x) = \frac{1}{117}$

\begin{solution}
$x = \arctan(117) + \pi k$, in  $[0, 2\pi)$, $x \approx 1.5622, \, 4.7038$
\end{solution}
\end{problem}

\begin{problem} \label{equarctriglast}
$\tan(x) = -0.6109$

\begin{solution}
$x =  \arctan(-0.6109) + \pi k$, in  $[0, 2\pi)$, $x \approx 2.5932, \, 5.7348$
\end{solution}
\end{problem}

\end{question}

\begin{question}
In Exercises \ref{rewritesinusoidfirst} - \ref{rewritesinusoidlast}, rewrite the given function as a sinusoid of the form $C(t) = A \cos(\omega t + \phi)$ and $S(t) = A\sin(\omega t + \phi)$ (See Example \ref{expandedsinusoidinverseex}.)  Approximate the value of $\phi$ (which is in radians, of course) to four decimal places.

\begin{problem} \label{rewritesinusoidfirst}
$f(t) = 5\sin(3t) + 12\cos(3t)$ 

\begin{solution}
$f(t) = 5\sin(3t) + 12\cos(3t) = 13\sin\left(3t + \arcsin\left(\frac{12}{13}\right)\right) \approx 13\sin(3t + 1.1760)$

$f(t) = 5\sin(3t) + 12\cos(3t) = 13\cos\left(3t + \arcsin\left(-\frac{5}{13}\right)\right) \approx 13\cos(3t -0.3948)$
\end{solution}
\end{problem}

\begin{problem}
$f(t) = 3\cos(2t) + 4\sin(2t)$

\begin{solution}
$f(t) = 3\cos(2t) + 4\sin(2t) = 5\sin\left(2t+\arcsin\left(\frac{3}{5}\right) \right) \approx 5\sin(2t+0.6435)$

$f(t) = 3\cos(2t) + 4\sin(2t) = 5\cos\left(2t+\arcsin\left(-\frac{4}{5}\right) \right) \approx 5\cos(2t-0.9273)$
\end{solution}
\end{problem}

\begin{problem}
$f(t) = \cos(t) - 3\sin(t)$

\begin{solution}
$f(t) = \cos(t) - 3\sin(t) = \sqrt{10} \sin\left(t + \arccos\left(-\frac{3\sqrt{10}}{10} \right)\right) \approx \sqrt{10} \sin(t + 2.8198)$ \\

$f(t) = \cos(t) - 3\sin(t) = \sqrt{10} \cos\left(t + \arcsin\left(\frac{3\sqrt{10}}{10} \right)\right) \approx \sqrt{10} \cos(t + 1.2490)$
\end{solution}
\end{problem}

\begin{problem}
$f(t) = 7\sin(10t) - 24\cos(10t)$

\begin{solution}
$f(t) = 7\sin(10t) - 24\cos(10t) = 25\sin\left( 10t + \arcsin\left(-\frac{24}{25}\right)\right) \approx 25 \sin(10t-1.2870)$ \\

$f(t) = 7\sin(10t) - 24\cos(10t) = 25\cos\left( 10t + \pi + \arcsin\left(\frac{7}{25}\right)\right) \approx 25 \cos(10t+3.4254)$
\end{solution}
\end{problem}

\begin{problem}
$f(t) = -\cos(t) - 2\sqrt{2} \sin(t)$

\begin{solution}
$f(t) = -\cos(t) - 2\sqrt{2} \sin(t) = 3\sin\left(t+\pi + \arcsin\left(\frac{1}{3}\right)\right) \approx 3\sin(t+3.4814)$ \\


$f(t) = -\cos(t) - 2\sqrt{2} \sin(t) = 3\cos\left(t+ \arccos\left(-\frac{1}{3}\right)\right) \approx 3\sin(t+1.9106)$
\end{solution}
\end{problem}

\begin{problem} \label{rewritesinusoidlast}
$f(t) = 2\sin(t) - \cos(t)$

\begin{solution}
$f(t) = 2\sin(t) - \cos(t) = \sqrt{5}\sin\left(t  + \arcsin\left(-\frac{\sqrt{5}}{5}\right)\right) \approx \sqrt{5}\sin(t -0.4636)$ \\

$f(t) = 2\sin(t) - \cos(t) = \sqrt{5}\cos\left(t  + \pi + \arcsin\left(\frac{2\sqrt{5}}{5}\right)\right) \approx \sqrt{5}\cos(t + 4.2487)$
\end{solution}
\end{problem}

\end{question}

\begin{question}

In Exercises \ref{domainexerfirst} - \ref{domainexerlast}, find the domain of the given function.  Write your answers in interval notation.

\begin{problem} \label{domainexerfirst}
$f(x) = \arcsin(5x)$

\begin{solution}
$\left[-\frac{1}{5}, \frac{1}{5}\right]$
\end{solution}
\end{problem}

\begin{problem}
$f(x) = \arccos\left(\frac{3x-1}{2} \right)$

\begin{solution}
$\left[-\frac{1}{3}, 1 \right]$
\end{solution}
\end{problem}

\begin{problem}
$f(x) = \arcsin\left(2x^2\right)$ 

\begin{solution}
$\left[-\frac{\sqrt{2}}{2}, \frac{\sqrt{2}}{2}\right]$ 
\end{solution}
\end{problem}

\begin{problem}
$f(x) = \arccos\left(\frac{1}{x^2-4}\right)$

\begin{solution}
$(-\infty, -\sqrt{5}] \cup [-\sqrt{3}, \sqrt{3}] \cup [\sqrt{5}, \infty)$ 
\end{solution}
\end{problem}

\begin{problem}
$f(x) = \arctan(4x)$

\begin{solution}
$(-\infty, \infty)$
\end{solution}
\end{problem}

\begin{problem}  $f(x) = \text{arccot}\left(\frac{2x}{x^2-9}\right)$

\begin{solution}
$(-\infty, -3) \cup (-3,3) \cup (3, \infty)$
\end{solution}
\end{problem}

\begin{problem}
$f(x) =\arctan(\ln(2x-1))$

\begin{solution}
$\left(\frac{1}{2}, \infty \right)$
\end{solution}
\end{problem}

\begin{problem}
$f(x) = \text{arccot}(\sqrt{2x-1})$

\begin{solution}
$\left[\frac{1}{2}, \infty \right)$
\end{solution}
\end{problem}

\begin{problem}
$f(x) = \text{arcsec}(12x)$

\begin{solution}
$\left(-\infty, -\frac{1}{12}\right] \cup \left[\frac{1}{12}, \infty\right)$
\end{solution}
\end{problem}

\begin{problem}
$f(x) = \text{arccsc}(x+5)$

\begin{solution}
$(-\infty, -6] \cup [-4, \infty)$
\end{solution}
\end{problem}

\begin{problem}
$f(x) = \text{arcsec}\left(\frac{x^3}{8}\right)$

\begin{solution}
$(-\infty, -2] \cup [2, \infty)$
\end{solution}
\end{problem}

\begin{problem} \label{domainexerlast}
$f(x) = \text{arccsc}\left(e^{2x}\right)$ 

\begin{solution}
$[0, \infty)$
\end{solution}
\end{problem}

\end{question}

\begin{problem}
Find the following limits.

\begin{enumerate}

\item $\ds{ \lim_{x \rightarrow 1^{-}} \arcsin(x)}$

\begin{solution}
$\ds{ \lim_{x \rightarrow 1^{-}} \arcsin(x)}$ $=\frac{\pi}{2}$
\end{solution}

\item  $\ds{\lim_{x \rightarrow -\infty} \arctan(3x)}$

\begin{solution}
$\ds{\lim_{x \rightarrow -\infty} \arctan(3x) = -\infty}$
\end{solution}

\item  $\ds{\lim_{x \rightarrow 1} \text{arcsec}(2x)}$

\begin{solution}
$\ds{\lim_{x \rightarrow 1} \text{arcsec}(2x)}$ $= \frac{\pi}{3}$
\end{solution}

\end{enumerate}

\end{problem}

\begin{problem}
Find a nonzero number $x$ where $\mbox{arccot}(x) \neq \arctan \left( \frac{1}{x} \right)$.
\end{problem}

\begin{problem}
Find an example where $\mbox{arcsec}(x) \neq \arccos \left( \frac{1}{x} \right)$ if we use $\left[0, \frac{\pi}{2} \right) \cup \left[ \pi,  \frac{3\pi}{2} \right)$ as the range of $f(x) = \mbox{arcsec}(x)$.
\end{problem}

\begin{problem}
Show that $\arcsin(x) + \arccos(x) = \frac{\pi}{2}$ for $-1 \leq x \leq 1$.
\end{problem}

\begin{problem}
Discuss with your classmates why $\arcsin\left(\frac{1}{2}\right) \neq 30$.
\end{problem}



\begin{problem}
Use the diagram below along with the accompanying questions to show: \[\arctan(1) + \arctan(2) + \arctan(3) = \pi\]

\begin{center}

\begin{image}
\begin{tikzpicture}[scale=1,font=\small]
\begin{axis}[
    xmin=-1, xmax=2.25,
    ymin=-1, ymax=3.25,
    width=3in,
    height=3in,
    axis x line=center,
    axis y line=center,
    xtick={0,1,2},
    ytick={1,2,3},
    xticklabels={$0$,$1$,$2$},
    yticklabels={$1$,$2$,$3$},
    tick style={line width=1.2pt},
    major tick length=6pt,
    axis line style={line width=1.2pt},
    x axis line style={-{Latex[length=10pt,width=8pt]}},
    y axis line style={-{Latex[length=10pt,width=8pt]}},
    xlabel={$x$},
    ylabel={$y$},
    every axis x label/.style={
        at={(current axis.right of origin)},
        anchor=west
    },
    every axis y label/.style={
        at={(current axis.above origin)},
        anchor=south
    }
]

% Line segments
\addplot[black,line width=2pt] coordinates {(0,1) (1,0)};
\addplot[black,line width=2pt] coordinates {(1,0) (2,3)};
\addplot[black,line width=2pt] coordinates {(0,1) (2,3)};
\addplot[black,line width=2pt] coordinates {(2,0) (2,3)};

% Points
\addplot[
    only marks,
    mark=*,
    mark size=3pt
]
coordinates {
    (0,0)
    (1,0)
    (2,0)
    (2,3)
    (0,1)
};

% Labels
\node[anchor=east]  at (axis cs:0,1) {$A(0,1)$};
\node[anchor=north east] at (axis cs:0,0) {$O(0,0)$};
\node[anchor=south] at (axis cs:1,-0.63) {$B(1,0)$};
\node[anchor=south] at (axis cs:2,-0.63) {$C(2,0)$};
\node[anchor=west] at (axis cs:1.25,3) {$D(2,3)$};

% Angle labels
\node at (axis cs:0.63,0.12) {$\alpha$};
\node at (axis cs:0.88,0.27) {$\beta$};
\node at (axis cs:1.18,0.14) {$\gamma$};

\end{axis}
\end{tikzpicture}
\end{image}

\end{center}

\begin{enumerate}

\item Clearly $\triangle AOB$ and $\triangle BCD$ are right triangles because the line through $O$ and $A$ and the line through $C$ and $D$ are perpendicular to the $x$-axis.  Use the distance formula to show that $\triangle BAD$ is also a right triangle (with $\angle BAD$ being the right angle) by showing that the sides of the triangle satisfy the Pythagorean Theorem.

\item Use $\triangle AOB$ to show that $\alpha = \arctan(1)$
\item Use $\triangle BAD$ to show that $\beta = \arctan(2)$
\item Use $\triangle BCD$ to show that $\gamma = \arctan(3)$

\item Use the fact that $O$, $B$ and $C$ all lie on the $x$-axis to conclude that $\alpha + \beta + \gamma = \pi$.  Thus $\arctan(1) + \arctan(2) + \arctan(3) = \pi$.

\end{enumerate}

\end{problem}

\end{document}

